\chapter{The Ultimate Choice}

This lecture proceeds by widening and then testing our moral field of view. Sandel begins with March 11 as a question for the world, not only for Japan; from there he lets the discussion move from solidarity and sacrifice, to the unequal distribution of nuclear risk, to the scale of cross-border harm, and finally to the question whether human sympathy can genuinely become global. The formal structure is conceptual rather than algebraic, but it is still precise: each new step in the lecture sharpens a relation, introduces a contrast, or tests an analogy.

\section{A Global Forum After March 11}

We begin not with a case in the narrow sense but with a frame. Sandel presents the earthquake, tsunami, and nuclear crisis as the subject of a special lecture on how we should live after March 11 and what this disaster reveals about human values. That opening matters because it sets the scale of everything that follows. Before we ask about nuclear power, or sympathy, or citizenship, we are asked to see the event as morally and politically world-sized.

The lecture's first formal move is therefore:
\begin{equation}
\text{March 11 disaster}
\Longrightarrow
\text{ethical question of how to live after it}.
\end{equation}
This is not a conclusion drawn at the end; it is the premise with which the whole discussion begins. The disaster is not merely an object of reporting. It is already being treated as an occasion for judgment.

\section{Solidarity, Sacrifice, and the National Family}

The first substantive contribution identifies a striking feature of the Japanese response. People rely on one another, they are willing to sacrifice for one another, and they continue to do what is needed even at the nuclear plants for the sake of fellow citizens. Sandel lets this stand before any theoretical dispute begins. The lecture first wants us to register a form of conduct.

We can write the relation this way:
\begin{equation}
\text{mutual reliance} + \text{sacrifice}
\Longrightarrow
\text{social solidarity}.
\end{equation}
A Japanese participant then gives the same thought a more intimate form by saying that the country feels like one family:
\begin{equation}
\text{nation}
\approx
\text{one family}.
\end{equation}
The approximation sign is important. The lecture is not defining a nation as a family in any strict sense. It is recording the moral felt sense of closeness that disaster can produce.

\subsection{Question \& Answer}

\paragraph{Question.} How can disaster turn strangers into something like a moral family?

\paragraph{Answer.} The lecture's answer begins from practice, not theory. Under pressure, people rely on one another, sacrifice for one another, and act for fellow citizens. The family image is Sandel's participants' way of naming the intensity of that solidarity.

\section{Nuclear Power, Distributed Risk, and the Airplane Analogy}

The lecture now turns. The same discussion that celebrates solidarity becomes a discussion of justice. A Japanese participant says that one problem with nuclear power is that the place bearing the risk is not the place receiving the benefit. This is the first clear asymmetry in the lecture:
\begin{equation}
\text{risk-bearing location}
\neq
\text{benefit-receiving location}.
\end{equation}
That relation changes the level of the discussion. A society may feel united, but institutions can still distribute burdens unfairly within it.

Sandel then lets the strongest everyday defense of dangerous technology appear. Airplanes sometimes break down, the consequences can be terrible, and yet we do not stop using them. Sometimes they are the only means to the end. In compact form:
\begin{align}
\text{airplane disasters}
&\not\Rightarrow
\text{stop using airplanes},\\
\text{airplanes are sometimes the only means to the end}
&\Longrightarrow
\text{continued use despite disaster}.
\end{align}
Only then comes the analogical leap:
\begin{equation}
\text{nuclear energy}
\stackrel{?}{\sim}
\text{airplane technology}.
\end{equation}
The question mark matters. The lecture does not accept the analogy merely because it has been proposed. It pauses over it and asks whether the analogy really holds once we take the structure of nuclear risk seriously.

\subsection{Question \& Answer}

\paragraph{Question.} When may we keep using a dangerous technology whose risks and benefits fall on different people?

\paragraph{Answer.} The first answer given in the lecture is pragmatic: we sometimes continue to use dangerous technologies when they are necessary means to important ends. But the lecture does not let that answer stand by itself, because the distribution of risk and benefit has already put pressure on the analogy.

\paragraph{Worked derivation.} The argumentative structure here is worth laying out step by step:
\begin{enumerate}
\item Nuclear power imposes risks on one location while delivering benefits elsewhere.
\item Dangerous technologies are not always abandoned after catastrophe.
\item Airplanes provide the comparison case.
\item Therefore the defense of nuclear power depends on whether nuclear danger is morally analogous to ordinary technological danger.
\end{enumerate}
This is the central worked derivation of the chapter's first half. The lecture is not merely contrasting opinions; it is building a test for a proposed analogy.

\section{Scope, Scale, and a Global Crisis}

The analogy is then resisted, not by denying that airplanes can fail, but by insisting that nuclear power differs in scope and scale. This is where the lecture moves from unequal distribution within a nation to consequences beyond the nation:
\begin{equation}
\text{scope and scale of nuclear crisis}
>
\text{scope and scale of airplane breakdown}.
\end{equation}
The point is then made concrete:
\begin{equation}
\text{nuclear leak in Japan}
\Longrightarrow
\text{effects in China and America}.
\end{equation}
And from this follows the widened political conclusion:
\begin{equation}
\text{global effects}
\Longrightarrow
\text{world's attention and world's effort}.
\end{equation}
This is one of Sandel's characteristic enlargements. What began as a question about a dangerous technology and an unjust distribution of risk now becomes a question about a harm whose effects spill across borders. The lecture has moved from domestic justice to transnational consequence.

\subsection{Question \& Answer}

\paragraph{Question.} What makes nuclear risk morally different from other dangerous technologies?

\paragraph{Answer.} The lecture's answer is not that nuclear power is uniquely risky in every respect, but that the scope and scale of its failure are different. Once the consequences spread far beyond the place of origin, the moral and political response can no longer be contained within the ordinary frame of local technological risk.

\section{Rousseau and the Problem of Distance}

Only after the lecture has widened the crisis does Sandel introduce a philosophical obstacle. He quotes Rousseau, who says that the sentiment of humanity evaporates and weakens as it is extended over the entire world. In the lecture's formal idiom we can compress this as
\begin{equation}
\text{sentiment of humanity}
\downarrow
\quad \text{as it is extended over the entire world}.
\end{equation}
Rousseau is being used here as the philosopher of moral distance. It may be easy to care for those near us; it may be much harder to care in an equally vivid way for humanity in general. Sandel brings in Rousseau precisely when the discussion has become global, because now the question is no longer only whether a disaster has global effects, but whether our sympathy can genuinely match that scale.

\subsection{Question \& Answer}

\paragraph{Question.} Can sympathy remain strong when its object is distant and global?

\paragraph{Answer.} Rousseau's answer is no: concern weakens as it spreads. Sandel uses that answer as a serious challenge, not as a decorative quotation. The lecture now has to ask whether a global moral horizon is psychologically and politically sustainable.

\section{Media Nearness, Community, and the Limits of Global Citizenship}

The lecture's reply to Rousseau comes through the modern experience of mediated nearness. If Rousseau were alive today, one participant suggests, and could watch tsunami footage online, he might no longer regard distant suffering as something at the far edge of the world. The lecture's counter-thought is therefore:
\begin{equation}
\text{geographical distance}
\not\Rightarrow
\text{moral distance}.
\end{equation}
Another participant adds that communication matters, and that in the case of a natural disaster we are brought together as a community:
\begin{equation}
\text{natural disaster}
\Longrightarrow
\text{community}.
\end{equation}
This is a strong answer to Rousseau, but Sandel does not let it harden too quickly into a doctrine of universal citizenship. One speaker says we can sympathize with countries half a world away; another remains skeptical that this amounts to a genuine global civic identity. The distinction should therefore remain explicit:
\begin{equation}
\text{global sympathy}
\neq
\text{global citizenship}.
\end{equation}
This is exactly the kind of distinction the lecture wants to preserve. Media can collapse distance; communication can widen concern; a disaster can make us feel part of a larger human community. But none of this automatically yields a full cosmopolitan identity.

\subsection{Question \& Answer}

\paragraph{Question.} Does global sympathy imply global citizenship, or only a thinner human solidarity?

\paragraph{Answer.} The lecture's answer is deliberately mixed. It allows that communication can make suffering half a world away morally vivid, but it stops short of claiming that such sympathy amounts to a fully shared civic membership. The widening of concern is real; the completion of a global identity remains doubtful.

\section{Human Pride and the Closing Horizon}

The lecture closes not with a final doctrine but with an emotion that is morally informative. A participant says she felt pride in the fact that people were not looting or hoarding, and in learning of the actions of Japanese people who acted as heroes. This final sentiment can be written as
\begin{equation}
\text{human pride}
\Longrightarrow
\text{admiration for non-looting, non-hoarding, and heroism}.
\end{equation}
That closing matters because it gives the lecture its final scale. We do not end with a proof of cosmopolitan citizenship, and we do not collapse back into narrow national feeling. We end with a widened moral sentiment: a pride in human conduct under disaster, broader than patriotism but thinner than a universal political identity.

Sandel's final hope matches that register. He wants the discussion to help both the people of Japan and those around the world who wish to support them find something in the argument that clarifies what this disaster has demanded of us.

\section{Summary}

The lecture unfolds by measured enlargements. It begins by treating March 11 as an ethical question for an international public,
\begin{equation}
\text{March 11 disaster}
\Longrightarrow
\text{ethical question of how to live after it},
\end{equation}
then identifies solidarity through mutual reliance and sacrifice,
\begin{equation}
\text{mutual reliance} + \text{sacrifice}
\Longrightarrow
\text{social solidarity},
\end{equation}
then turns that solidarity into a problem of justice through the asymmetry
\begin{equation}
\text{risk-bearing location}
\neq
\text{benefit-receiving location}.
\end{equation}
From there Sandel stages the airplane analogy, resists it by appeal to scope and scale, globalizes the consequences of the crisis, and then asks whether human sympathy can widen in step with those consequences. Rousseau articulates the problem of distance; communication and media offer a reply; but the lecture closes by distinguishing global sympathy from global citizenship and by ending, more cautiously, in human pride.

What must be preserved in the final chapter is this rhythm of pressure and response. We begin with solidarity, then test it against institutional asymmetry; we widen the crisis, then test whether our moral imagination can widen with it. The lecture does not simply list positions. It moves from one sharpened relation to the next, each one motivated by the limits of the one before.